\subsection{Motivation}

It started with curiosity. I found a piece of Danish text on the Sverdrup Society page and wondered if it could be translated by AI. This is what came back after several hours of prompt forming:

\begin{quote}

The formation of the Church is the holy work of the triune God. 

Sin had marred creation, and death ruled over the fallen world; divine wrath rested upon humanity, once created for salvation. Yet God’s eternal plan—conceived in love and grace—sought to restore what sin had corrupted, to overthrow the powers of sin, death, and condemnation, and to grant new life and immortality to those enslaved by death. Therefore, no congregation could come into being, no fellowship of God’s redeemed children could be established, except through reconciliation in Christ’s death, through life raised from the dead in his resurrection, and through union with God by the Holy Spirit. God’s free and living people on earth share their closest communion with him through the death and resurrection of his Son and the fellowship of the Holy Spirit.

\end{quote}


Some time later this emerged from the machine translation: 

\begin{quote}
All hands on deck, draw the sword from the sheath, let the cannons glint in the sun, strike upon the souls, that the sound of weapons may be heard far throughout the land. The time has come, the air is pregnant with thunder; let the thunder roll in the cloud and the lightnings strike with flash upon flash in the murk of night. The powers of the world are stirred; let the powers of heaven stir themselves in us. God’s Church is in danger; let weapon-bearing men defend it. Against the lie we use sharp swords, against the lie we bear flashing helmets, against the lie we take up bright shields. \textbf{But the lie is not outside God’s Church, but within its walls.} Therefore awake, you Christian Church, wherever you are found; therefore awake, you watchman of Zion, and blow the trumpet, that it may give a clear sound. 

The lie is this, that Christ is dead; for behold, he lives. 

The lie is this, that Christianity is mere doctrine and the understanding of understandings and the clarity of thought; for behold, it is the life of love and the folly of love. 

The lie is this, that the Church is an institution for ruling over men; for behold, it is the home of freedom and the fellowship of the liberated. 

The lie is this, that the Church must use the world’s falseness and deceit in order to be shrewd as a serpent; for behold, it is the holy witness of truth and simplicity. 

The lie is this, that the world is the Church’s most dangerous enemy; for behold, it is the object of God’s love. 

The lie is this, that the Church is pure; for behold, it is full of rottenness and stench. 

The lie is this, that Christianity is hatred of the world; for behold, it is love for the world. 

The lie is this, that the Church must fear the world; for behold, the world has been overcome. 

The lie is this, that the Church must go outward to find its life; for behold, its enemy is in the midst of it. Let us find him where his work is among us; let us know him as the one he is. Let us meet him as Christ did of old: I am the one you seek.

\end{quote}

Wow, could Sverdrup write! I was hooked. From that point forward, I spent hundreds of hours machine translating Sverdrup. This included the entire book of sermons Aand og Liv (Spirit and Life) by Professors Oftedal and Sverdrup. 

Patterns started to emerge. The first was a preaching style where Sverdrup would take hold of you with an old testament lessons until you knew—deep down, you really knew—that you were the old Adam. And then Sverdrup would swing the hinge revealing the freedom of the Gospel. 

But Sverdrup was not merely a preacher he was a professor to preachers. There we see the deeper patterns where he serves as the defender of the congregation. Sverdrup has clarity on this topic with the Fundamental Principles at the top of a large body of work. These are the same principles the AFLC holds as presented in the front matter of this book.


\subsection{Evjen}

Up until this point, Sverdrup had been a fascinating project to which I had dedicated half a year's free time translating and reading many hundreds of pages. The work was serious, interesting, and enjoyable, but I had yet to find what I was looking for. Specifically, how is a person to act beyond the congregation: if the congregation is a vessel, how should the sailor act when ashore?

At a surface level I was attempting to reconcile Sverdrup both Martin Luther and Gene Edward Veith. I knew there was a relationship but I had no way to describe it. I had the pleasure of exploring these ideas with Loiell Dyrud who many of you know. 

Almost as an afterthought, Loiell placed an old faded copy of a Norwegian book called Guidance into my hands. This book contained material from Sverdrup and Oftedal compiled by Evjen - see section 3. A month later, Professor Martin Horn asked me to use the machine to translate Evjen's The Position - see section 4. Ironically, it was Evjen that allowed me to clearly see Sverdrup. There was one night where the entire structure snapped into position and I was able to see a cohesive set of rules starting from the Pauline language in Ephesians, through Luther, through Sverdrup, and yes even through Veith.

That nautical understanding is developed in the next two sections of this book where the congregation is compared to a working vessel with all hands on deck. To be clear, Evjen forced me to critically examine what Sverdrup was saying. Veith provides the modern language for what Sverdrup calls spirit and life and Paul calls the body of Christ. To be alive is more than doctrine. It is to be the hands of God, embodied in faithful vocation, in the service of daily bread for family, neighbor, community, and nation.


  
     

   



   
   
  